% ch3_c §3.4–3.5 (书页 91–97 / p105–p111)
%--B-TAIL--
%JOIN%速度。例如，在简单的区边界附近，至少必须把两支波组合起来，得到一个类似于式 (3.17) 的二次方程，它给出电子或光子能量作为 $\mathbf{k}$ 的函数的两个根。

这些根生成了图 50 中的\emph{色散面}。原先以倒格点 $P$ 与 $Q$ 为中心、半径为 $\nu/c$ 的“自由光子”球面本应在区边界上的 $O$ 点相交（参看图 29 与图 37），如今则分裂成两个双曲面。因此，图中诸如 $S$ 或 $S'$ 这样的点代表一个频率为 $\nu$ 的电磁扰动，它以两支平面波的组合形式在晶体中传播，其波矢为 $SP$ 与 $SQ$，或 $S'P$ 与 $S'Q$。为了与这个频率的入射 X 射线相匹配，色散面两个分支的波束都被激发，从而导致 §2.6 中提到过的干涉效应。除了尺度上的差别——X 射线或快电子的波矢通常远大于倒格子的典型矢量——图 50 其实与图 42 完全相同。

\section{紧束缚法}

N.F.E. 方法着眼于原子芯以外的波函数，它们看上去非常像平面波；而在芯区内，它们则像原子轨道。这提示了一种截然不同的电子波函数构造方案：我们试着把各自局域在特定原子上的原子轨道组合起来，用以表示一个贯穿整个晶体的状态。

设 $\phi_a(\mathbf{r}-\mathbf{l})$ 是中心位于 $\mathbf{l}$ 处的自由原子的一个原子轨道。于是我们可以构造一个满足 Bloch 条件 (1.58) 的函数
\begin{equation*}
\phi_{\mathbf{k}}(\mathbf{r})=\sum_{l}e^{i\mathbf{k}\cdot\mathbf{l}}\phi_a(\mathbf{r}-\mathbf{l}).
\tag{3.24}
\end{equation*}
这个函数看起来像一列强局域化的原子轨道，乘上一个波动的相位因子 $\exp(i\mathbf{k}\cdot\mathbf{l})$。在每个原子内部，局域轨道占主导，应当是局域 Schr\"odinger 方程的一个很好的解。

作为一级近似，让我们对这个波函数计算能量的期望值
\begin{equation*}
\mathcal{E}(\mathbf{k})=\frac{\displaystyle\int\phi_{\mathbf{k}}^{*}\left\{-\frac{\hbar^{2}}{2m}\nabla^{2}+\mathcal{V}(\mathbf{r})\right\}\phi_{\mathbf{k}}\,\mathrm{d}\mathbf{r}}{\displaystyle\int\phi_{\mathbf{k}}^{*}\phi_{\mathbf{k}}\,\mathrm{d}\mathbf{r}}.
\tag{3.25}
\end{equation*}

%FIG{51}: 固体中的波函数

如果相邻晶胞中轨道的交叠不太大，我们或许可以把分母中的归一化因子取为 1。分子可以写成
\begin{align*}
\mathcal{E}(\mathbf{k})&\approx\sum_{ll'}e^{i\mathbf{k}\cdot(\mathbf{l}-\mathbf{l}')}\int\phi_a^{*}(\mathbf{r}-\mathbf{l}')\left\{-\frac{\hbar^{2}}{2m}\nabla^{2}+\mathcal{V}(\mathbf{r})\right\}\phi_a(\mathbf{r}-\mathbf{l})\,\mathrm{d}\mathbf{r}\\
&=\sum_{\mathbf{h}}e^{i\mathbf{k}\cdot\mathbf{h}}\mathcal{E}_{\mathbf{h}},
\tag{3.26}
\end{align*}
其中
\begin{equation*}
\mathcal{E}_{\mathbf{h}}=\frac{1}{v_c}\int\phi_a^{*}(\mathbf{r}+\mathbf{h})\left\{-\frac{\hbar^{2}}{2m}\nabla^{2}+\mathcal{V}(\mathbf{r})\right\}\phi_a(\mathbf{r})\,\mathrm{d}\mathbf{r}.
\tag{3.27}
\end{equation*}
我们利用晶格势的周期性，使各个积分只依赖于轨道所居格点的相对位置 $\mathbf{h}=\mathbf{l}-\mathbf{l}'$。

现在设 $\phi_a(\mathbf{r})$ 满足方程
\begin{equation*}
\left\{-\frac{\hbar^{2}}{2m}\nabla^{2}+v_a(\mathbf{r})\right\}\phi_a(\mathbf{r})=\mathcal{E}_a\phi_a(\mathbf{r}),
\tag{3.28}
\end{equation*}
其中 $v_a(\mathbf{r})$ 现在是位于 $\mathbf{r}=\mathbf{0}$ 处的孤立原子的势。

%FIG{52}: （a）等能线；（b）紧束缚法得到的沿立方轴的 $\mathcal{E}(k)$

如果我们略去各种交叠积分，便得到
\begin{equation*}
\left.\begin{aligned}
\mathcal{E}_{0}&\approx\mathcal{E}_a,\\
\mathcal{E}_{\mathbf{h}}&\approx\frac{1}{v_c}\int\phi_a^{*}(\mathbf{r}+\mathbf{h})\left\{\mathcal{V}(\mathbf{r})-v_a(\mathbf{r})\right\}\phi_a(\mathbf{r})\,\mathrm{d}\mathbf{r}.
\end{aligned}\right\}
\tag{3.29}
\end{equation*}
显然，$\phi_a(\mathbf{r}+\mathbf{h})$ 与 $\phi_a(\mathbf{r})$ 之间必须有某种交叠，$\mathcal{E}_{\mathbf{h}}$ 才不至于为零——而 $\mathbf{h}$ 处原子附近的高势会增强这种交叠。但是我们可以预期，除了近邻以外，$\mathcal{E}_{\mathbf{h}}$ 都应当非常小。

考虑一个简单立方晶格，取 $\mathbf{h}=(a,0,0)$ 等等。于是，当 $\phi_a(\mathbf{r})$ 具有立方对称性时，
\begin{equation*}
\mathcal{E}(\mathbf{k})\approx\mathcal{E}_a+2\mathcal{E}_{100}(\cos ak_x+\cos ak_y+\cos ak_z),
\tag{3.30}
\end{equation*}
这个情形下的区当然是一个立方体，能面如图 52 所示。由式 (3.29) 易见 $\mathcal{E}_{100}$ 本质上是负的。我们得到一条宽度为 $12|\mathcal{E}_{100}|$ 的能带，其最低能量为 $\mathcal{E}_a+6\mathcal{E}_{100}$。

沿立方轴方向我们看到，$\mathcal{E}(\mathbf{k})$ 在 $k=0$ 处有一个极小，而在区边界 $\mathbf{k}=(\pi/a,0,0)$ 处有一个极大。这与 N.F.E. 模型的曲线（图 36）表面上有几分相似，但带底的曲率依赖于 $\mathcal{E}_{100}$。整条能带的结构依赖于诸如式 (3.29) 那样的积分，它们对原子芯之外轨道的细节以及晶格间距都非常敏感。显然，如果原子轨道高度局域化，或者晶格间距变大，能带将趋于变得非常窄。

这种\emph{紧束缚法}演示了一条重要的原理。设想我们取 $N$ 个原子，并让它们彼此相距很远。每个原子上将有各不相同的原子能级 $\phi_a,\phi_b,\phi_c$；而在整个集合中，对单个电子来说，这些将是 $N$ 重简并的状态。

%FIG{53}: 原子彼此靠拢时，原子能级展宽成带

但当我们把原子聚到一起时，轨道发生交叠——于是我们发现，$N$ 重简并被劈裂成一个态带。每一条原子能级都将产生一个完整的、含 $N$ 个态的能带——因此我们可以谈及由相应原子态产生的 3s 带、4p 带等等。这种分类在过渡金属的情形下特别贴切：原子的 d 态在每个芯内都相当紧致，因而它们组合成一个相对较窄且界定清晰的 d 带。

然而，从图 53 可以清楚地看到，不同的能带可能展宽到彼此相交。这时就必须修改我们在此讨论的简单方案。我们必须假定波函数是\emph{原子轨道的线性组合}（L.C.A.O.）
\begin{equation*}
\psi_{\mathbf{k}}=\sum_{lj}e^{i\mathbf{k}\cdot\mathbf{l}}\beta_j\phi_a^{(j)}(\mathbf{r}-\mathbf{l}),
\tag{3.31}
\end{equation*}
其中 $\phi_a^{(j)}(\mathbf{r}-\mathbf{l})$ 是 $\mathbf{l}$ 处原子上一组互不相同的原子轨道之一。

然后我们必须把它用作试探波函数来估算能量，并通过适当选取系数 $\beta_j$ 而使能量取极小。显而易见，我们将得到一组组线性方程，其系数中含有诸如式 (3.26) 中 $\mathcal{E}_{\mathbf{h}}$ 那样的表达式。这些表达式又涉及如式 (3.29) 那样的积分，只是其中的两个原子轨道也许对应于两个不同格点上的不同原子能级。

L.C.A.O. 方法在量子化学中被广泛用于构造\emph{分子}波函数，但它用于固体中 Bloch 函数的定量计算并不真正令人满意。我们不仅在三中心积分和非正交基函数上遇到计算困难；而且，用诸如式 (3.31) 那样的表达式去表示半导体或金属中的价电子态，还存在着根本性的异议。

当我们构成固体时，每个原子都与其近邻的势相交叠。让我们以一种简单化的想法（图 54 与图 55）假定：各别的原子势可以叠置起来，在以原子核为中心的主势阱之间的间隙区域中形成一个较低而且恒定的势。显然，式 (3.31) 中所用的那些单个原子轨道已不复存在，因为它们此时将位于原子球之间势垒的能量之上。原则上，试图用全然不同的势 $v_a(\mathbf{r})$（它满足无关的边界条件：当 $r\to\infty$ 时 $\phi_a(\mathbf{r})$ 迅速趋于零）的本征态来表示势 $\mathcal{V}(\mathbf{r})$ 中的 Bloch 函数，并不是一个好方法。所有的节点都错置了位置，而且不容易挪动。

%FIG{54}: 自由原子的束缚原子轨道

%FIG{55}: 晶体的 Bloch 态

无论如何，这组函数是\emph{不完备}的，因为它缺少 Schr\"odinger 方程在连续谱中、高于 $v_a(\mathbf{r})$ 能量零点的所有散射波本征态。虽然式 (3.31) 在原子芯内部可以与 $\psi_{\mathbf{k}}$ 匹配得相当好，但它谈不上能表示间隙区域中的 Bloch 态——在那里它必须表现得像自由电子平面波的简单组合（见 §3.6）。

不过，这个方法确实给出了一条重要的形式原理。

在式 (3.26) 中我们看到，$\mathbf{k}$ 空间中的能量可以展开成 Fourier 级数
\begin{equation*}
\mathcal{E}(\mathbf{k})=\sum_{l}e^{i\mathbf{k}\cdot\mathbf{l}}\mathcal{E}_l.
\tag{3.32}
\end{equation*}
可以证明，这种形式的展开总是可能的——也就是说，L.C.A.O. 方法如果精确而彻底地进行下去，将会给出这种形式的能量——尽管系数 $\mathcal{E}_l$ 不再是像式 (3.27) 或 (3.29) 那样的简单积分，而且当 $\mathbf{l}$ 很大时也许收敛得不太快。

现在请记住，正格矢 $\mathbf{l}$ 本身就是倒格子 $\mathbf{g}$ 的倒格矢（§1.3 (iv)）。式 (3.32) 正是如下这类公式
\begin{equation*}
\mathcal{V}(\mathbf{r})=\sum_{\mathbf{g}}e^{i\mathbf{r}\cdot\mathbf{g}}\mathcal{V}_{\mathbf{g}}
\tag{3.33}
\end{equation*}
的对偶（译注：原书此式编号误印为“(3.23)”。）；只要 $\mathcal{V}(\mathbf{r})$ 具有周期性 $\mathcal{V}(\mathbf{r}+\mathbf{l})=\mathcal{V}(\mathbf{r})$，这类公式便成立。于是，\emph{每条能带中的能量函数在倒格子中必须是周期的}；
\begin{equation*}
\mathcal{E}(\mathbf{k}+\mathbf{g})=\mathcal{E}(\mathbf{k}).
\tag{3.34}
\end{equation*}

这个性质我们在上一节中已经注意到。(3.32) 这种表示式有时颇为方便；它使 $\mathcal{E}(\mathbf{k})$ 自动满足这个条件。如果把系数 $\mathcal{E}_l$ 当作可调参数来处理，往往就能用这个公式拟合观测到的费米面。作为内插公式，它的“物理”论证比 N.F.E. 方案略逊一筹——在后一方案中，Fourier 分量 $\mathcal{V}_{\mathbf{g}}$ 可以取作参数——但它更便于计算，因为它把能量直接表示成简单解析函数之和。

\section{元胞法}

L.C.A.O. 方法所受的非议之一，可以通过使用被构造得满足 Bloch 条件 (1.58) 的函数来避免；也就是说，当波函数移动通过晶格平移 $\mathbf{l}$ 时（例如从晶胞的一个面移动到另一个面），它改变一个因子 $\exp(i\mathbf{k}\cdot\mathbf{l})$。另一种做法是，对公式
\begin{equation*}
\psi_{\mathbf{k}}(\mathbf{r})=e^{i\mathbf{k}\cdot\mathbf{r}}u_{\mathbf{k}}(\mathbf{r}),
\tag{3.35}
\end{equation*}
中的函数 $u_{\mathbf{k}}(\mathbf{r})$ 建立微分方程，同时（由式 (1.61)）知道它必须是正格子中的周期函数。

在一维情形，这看上去是个简单的方案。例如，可以选取一个 $\mathcal{E}$ 值，积分方程，看条件是否满足，再改变 $\mathcal{E}$，如此等等。但在三维情形，它是笨拙的。条件 (3.34) 必须在晶胞的整个面上成立——这就意味着要在无穷多个点上成立。于是，人们在晶胞内构造各种球谐函数的解，并寻找在晶胞边界上有限个点处相匹配的线性组合。这还可以改进：在晶胞表面上定义一个平均匹配误差，并通过改变系数使其极小。这一步骤在数学上很繁琐，但原理上足够明了；它曾为某些固体给出过相当好的能带结构，但在新理论概念方面却少有贡献。

元胞法看上去简单而直接。晶胞是一个定义明确的几何概念，Bloch 条件也很容易写下。但是整个想法并不很对头。我们不应当人为地设立晶胞边界，偏偏在势最为平坦、波函数最接近自由电子波函数的区域制造匹配间断。匹配条件只能任意地定义，而有时表面上看去合理的做法竟给出完全错误的结果，甚至连空晶格这种初等情形也不例外。

这种一般性的元胞方法有一个前身，有时颇为有用。这就是 \emph{Wigner--Seitz 方法}。我们注意到，F.C.C. 或 B.C.C. 格子的晶胞——即 Wigner--Seitz 原胞——是一个近似于球体的规则多面体。姑且假定它是一个球，体积与晶胞相同，即半径为 $r_s$（参看式 (2.54)）。

考虑 $\mathbf{k}=0$ 的态。它必须满足
\begin{equation*}
\psi_{\mathbf{k}}(\mathbf{r}+\mathbf{l})=\psi_{\mathbf{k}}(\mathbf{r});
\tag{3.36}
\end{equation*}
即它必须逐个晶胞地周期。这意味着它在晶胞边界上应当有水平切线；我们可以说它必须满足条件
\begin{equation*}
\left.\frac{\partial\psi_{\mathbf{k}}}{\partial r}\right]_{r=r_s}=0
\tag{3.37}
\end{equation*}
于 Wigner--Seitz 球面上。这样一来，如果势在球内是球对称的，我们只需在这个边界条件下求解一个径向方程，于是能量 $\mathcal{E}(0)$ 便可唯一地确定（图 56）。

%FIG{56}: Wigner--Seitz 方法
